\zlabel{triptych:concise:themes:start}
\section{The Propers: Themes and Movement}

The Gospel can become fruitful among people who suffer, and it changes the desires and actions of those who receive it. The mustard seed and leaven describe a change whose outcome exceeds its beginning; Paul's thanksgiving names the faith, charity and endurance through which that change becomes visible.

\subsection{Promise, reflection and a community's example}

The Mass begins with God's promise of peace and gathering. In Jeremiah 29:11,12,14, that promise addresses exiles who still have a life to conduct away from home. Jerome explains that they must wait rather than follow promises of immediate return. God summons them to pray and seek him wholeheartedly. The Introit's psalm verse, Psalm 84(85):2, recalls favor already received. Its full psalm asks for mercy again. Gratitude and renewed petition can therefore belong to the same people.

The Collect asks that reflection become speech and action acceptable to God. Its historical English describes the object of reflection as what is reasonable and just. Schuster emphasizes this passage from attention to conduct. In the Epistle, Paul gives thanks for people whose actions he can describe: they have received the word under affliction, turned from idols and become an example to believers in Macedonia and Achaia. Their conversion has become known beyond the community that first received the preaching.\footnote{Jerome, \textit{Commentary on Jeremiah}, V, on 29:10--14, PL 24, cols.~893--894; Schuster, \textit{The Sacramentary}, I, p.~416; 1 Thessalonians 1:2--10.}

Aquinas explains how this example comes about. Faith must issue in action and endure its demands; charity abounds in works; hope sustains patience in adversity. Chrysostom especially develops the cost of continuing to love when opposition threatens that attachment. Their explanations connect belief with what other people experience of a believer. The Collect's words and deeds are consequently pertinent to both interpretations developed here: a community communicates its faith through conduct, while each person's understanding must become willing action.

\subsection{Praise and a cry for help}

The Gradual praises God for deliverance in Psalm 43(44):8--9. The full psalm soon describes defeat and suffering among people who maintain that they have kept the covenant. Bellarmine first hears the selected words as remembrance of actual rescues; Augustine also hears the certainty of future deliverance expressed as already accomplished. Present hardship remains intelligible in both explanations.

Before the Gospel, the Alleluia asks to be heard from the depths. Augustine understands the opening of Psalm 129(130) through the sinner's need for mercy. The person who recognizes that need can cry to God, whereas prosperity may conceal it. The Epistle's joy amid affliction and the psalmist's urgent prayer describe different aspects of dependence on God: the Spirit sustains believers in their suffering, and those believers still need to ask for help.\footnote{Aquinas, \textit{Commentary on First Thessalonians}, 1, lecture 1 [87862], on 1:3--6; Chrysostom, \textit{Homilies on First Thessalonians}, homily 1; Bellarmine, \textit{Commentary on the Book of Psalms}, Psalm 43, explanations 7--9; Augustine, \textit{Expositions of the Psalms}, 43.6--8 and 129.1--2.}

\clearpage
\subsection{The seed received and the leaven at work}

The Gospel tells of a seed becoming a tree and leaven hidden in meal until the whole is leavened. Chrysostom attends to their place after the sower's losses and the enemy's weeds. The disciples could fear that few people would ever believe. Christ answers with images of growth through small beginnings. For Chrysostom, the three measures mean a large quantity; the little leaven affects much meal through contact with it. He urges disciples to remain among the people to whom they are sent.

Jerome attends to the person receiving the preaching. The announcement of a crucified God may seem unimpressive beside eloquent philosophy. Received and nourished by faith, it grows into understanding that directs a life. He interprets the three measures as spirit, soul and body, or alternatively as reason, anger and desire. Scriptural understanding brings these into agreement: reason attains prudence, anger opposes vice, and desire seeks virtue. Chrysostom explains the reach of the Gospel among many; Jerome explains the completeness of its effect within the hearer.\footnote{Chrysostom, \textit{Homilies on Matthew}, 46.2; Jerome, \textit{Commentary on Matthew}, II, on 13:31--33, PL 26, cols.~89--92.}

Augustine identifies the leaven with love that warms and rouses the person; heart, soul and mind are one of his explanations of the three measures. Understanding must become a willing pursuit of the good. The appointed passage ends with Jesus speaking in parables and Matthew quoting Psalm 77(78):2. Chrysostom sees an invitation to inquiry; Jerome recalls the mysteries disclosed through Israel's history. The disciples continue learning by asking for an explanation.\footnote{Augustine, \textit{Questions on the Gospels}, I.12.}

\subsection{Offering, thanksgiving and continued desire}

At the Offertory, the cry from the depths returns with its opening repeated. Schuster describes the gifts as already received from God's goodness. The Secret asks that the offering cleanse, renew, govern and protect the people who bring it. For a community becoming an example, these petitions concern the conduct others will see. For a person learning to desire rightly, they concern forgiveness, renewed life and direction in choices still to come.

The Trinity Preface thanks the Father and confesses the Son and Holy Spirit with him as one God. The people whom Paul describes as serving the living and true God give that service definite expression in worship. Jerome reports a Trinitarian interpretation of the three measures, but warns that uncertain explanations of parables cannot establish dogma. The Preface's confession rests on God's revelation.

The Communion promises a hearing to believing prayer. Mark's following command requires forgiveness, and Bede directs those seeking pardon and eternal life to forgive their neighbors. After Communion, those already nourished ask for continued desire for the heavenly gifts. Schuster notices this continuing hunger after reception. Such desire concerns both perseverance in faithful witness and the reform of the person who must live it.\footnote{Chrysostom, \textit{Homilies on Matthew}, 47.1; Jerome, \textit{Commentary on Matthew}, II, on 13:33--35, PL 26, cols.~91--93; Schuster, \textit{The Sacramentary}, I, p.~417, III, p.~188; Mark 11:24--25; Bede, \textit{Commentary on Mark}, on 11:22--25.}
\zlabel{triptych:concise:themes:end}
